% NOETHER PAPER 2 — KOREAN TRANSLATION-PRODUCER DRAFT — UNCHECKED
% Unit: T19-U126; current-authority whole lines 2574--2587
% Source slice: 851 LF UTF-8 bytes; SHA-256 EEFB59F0E35E0B9EF31A099F0663EE59D3897F796C004D096EA4050B11344CCF
% Pointer: NOETH-DE-AUTH-v005-20260804; SHA-256 42E6844BFCBFB2133E9AA323A823604351CF9C49550AFCF34ECAAF7887185660
% This is producer translation only: no source/Korean/formula review, compilation, or rendering.

생기는 계를 다음 네 부분계로 나눈다.
\[
\begin{array}{ll}
\text{계 }s_I:&s\text{ 및 }s\text{의 거듭제곱들을 계 I과 수축하여 생긴다,}\\
\text{계 }s_{II}:&s\text{를 계 II와 수축하여 생긴다,}\\
\text{계 }s_{III}:&s\text{를 계 III의 개별 형식들(곱은 제외)과 수축하거나,}\\
&\text{계 I과 II에서 각각 한 형식씩 취한 곱과 수축하여 생긴다,}\\
\text{계 }s_{IV}:&s\text{를 계 I과 III, II와 III, III과 III에서 각각 한 형식씩}\\
&\text{취한 곱과 수축하여 생긴다.}
\end{array}
\]
또한 이제부터 모든 공식은 $(\varrho,t)$를 법으로 이해하고, 공식 (27.)부터는 $(\varrho,t,i,J,\text{ 더 높은 형식들})$을 법으로 이해한다. 이를 매번 명시하지는 않는다.

